% Next-steps plan for the supervisor, 2026-10-17.
% Build from this directory:  tectonic 2026-10-17-next-steps.tex
% Figures are the PNGs in graphs/ at the repo root (exported from notebook 14).
\documentclass[10pt,a4paper]{article}

\usepackage[a4paper,margin=2cm]{geometry}
\usepackage{graphicx}
\usepackage{xcolor}
\usepackage{caption}
\usepackage{enumitem}
\usepackage{titlesec}
\usepackage{float}
\usepackage{booktabs}
\usepackage{tabularx}
\usepackage[hidelinks]{hyperref}

\graphicspath{{../../graphs/}}

\definecolor{ink2}{HTML}{52514E}
\captionsetup{font={small,color=ink2},labelfont={small,color=ink2},justification=raggedright,
              singlelinecheck=false,skip=3pt}
\titleformat{\section}{\large\bfseries}{\thesection.}{0.5em}{}
\titlespacing*{\section}{0pt}{1.2em}{0.4em}
\setlist{itemsep=0.25em,topsep=0.3em}
\setlength{\parindent}{0pt}
\setlength{\parskip}{0.5em}
\setlength{\intextsep}{8pt}
\renewcommand{\arraystretch}{1.25}
\newcolumntype{L}{>{\raggedright\arraybackslash}X}

\hypersetup{pdftitle={Supplement KG next steps},pdfauthor={Yanfeng Li}}

% One full-width figure from graphs/.
\newcommand{\fig}[2]{%
  \begin{figure}[H]
    \centering
    \includegraphics[width=\linewidth]{#1}
    \caption{#2}
  \end{figure}}

\newcommand{\rel}[1]{\emph{#1}}

\begin{document}

{\LARGE\bfseries Supplement knowledge graph: observations and next steps}\\[0.3em]
{\small\color{ink2} Yanfeng Li \textperiodcentered{} 17 October 2026}

\section*{Summary}
\begin{itemize}
  \item iDISK 2.0's ingredient relations (\rel{is effective for}, \rel{has adverse reaction},
        \rel{interacts with}) all come from \textbf{MSKCC}, and cover a few hundred ingredients.
  \item DSLD and NHP contribute no relations. iDISK takes only names and composition from DSLD, and
        keeps NHP's Purpose and Risk as free text on the product.
  \item The extracted DSLD claims, NHP's Purpose/Risk texts and the \textbf{EU health-claims register}
        could all add ingredient relations. Doing so raises several open problems (Section~4).
\end{itemize}

\section{Observation: iDISK's relations come only from MSKCC}

\fig{05_idisk_relation_gap.png}{Effect edges in the iDISK 2.0 release, by source.}

\begin{table}[H]
\small
\begin{tabularx}{\linewidth}{@{}l r l L@{}}
\toprule
Relation (iDISK 2.0 release) & Edges & Source & Ingredients covered \\
\midrule
\rel{is effective for} (ingredient $\rightarrow$ disease)    & 827   & MSKCC only & 266 \\
\rel{has adverse reaction} (ingredient $\rightarrow$ symptom) & 1,342 & MSKCC only & 219 \\
\rel{interacts with} (ingredient $\rightarrow$ drug)        & 535   & MSKCC only & 179 \\
\rel{has ingredient} (product $\rightarrow$ ingredient)      & 317,062 & NHP (99.7\%), DSLD & -- \\
\bottomrule
\end{tabularx}
\end{table}

Only 285 of the 7,876 ingredients in iDISK have an effect edge, all from MSKCC's 279 herb
monographs. In iDISK 1.0, NMCD supplied most of these relations (e.g.\ 4,307 \rel{is effective for}
edges), but the 2.0 release files contain only the MSKCC edges.

\section{Observation: DSLD and NHP claims are not used as relations}

\begin{table}[H]
\small
\begin{tabularx}{\linewidth}{@{}l L L@{}}
\toprule
& \textbf{What iDISK 2.0 takes} & \textbf{What is left unused} \\
\midrule
DSLD (US labels) &
  Product name, company, ingredient names, \rel{has ingredient} links (but only 520 of the 92,442
  DSLD-only products have any). &
  All label statements. Purpose/Risk are filled only for the 208 products that are also in NHP
  (LNHPD). \\
NHP (Canadian licences) &
  Product name, ingredients, \rel{has ingredient}, and Purpose (65,097 products) and Risk (63,623)
  as \textbf{free text on the product}. &
  Which ingredient a Purpose or Risk is about. They never become relations, even for the 32,968
  products with a Purpose and exactly one ingredient. \\
\bottomrule
\end{tabularx}
\end{table}

iDISK builds its data from fields each source already structures (names, ingredient lists), and does
not extract anything from free text. Our extraction reads the label text itself: 466,630 purpose and
355,134 risk items from all 214,780 DSLD labels. So the claims now exist, but as text attached to a
product.

\section{Observation: the EU health-claims register}

The EU Register of nutrition and health claims lists, for each substance, the claims EFSA assessed and
its verdict: 269 \textbf{authorised} claims over 94 substances, and 2,067 \textbf{non-authorised}
claims over 1,177 substances.

\begin{table}[H]
\small
\begin{tabularx}{\linewidth}{@{}l r L@{}}
\toprule
EU claim type & Authorised & Example \\
\midrule
Art.\ 13: general function & 242 & ``Selenium contributes to the normal function of the immune system'' \\
Art.\ 14(1)(a): disease-risk reduction & 15 & ``Plant sterols have been shown to lower blood
  cholesterol. High cholesterol is a risk factor in the development of coronary heart disease.'' \\
Art.\ 14(1)(b): children's development & 12 & Function claims restricted to children. \\
\bottomrule
\end{tabularx}
\end{table}

Compared with label claims, an EU claim already names its \textbf{substance}, not a product, and
carries a \textbf{regulator's verdict}. Its wording is close to label wording: authorised claims read
like NHP Purpose, and rejected ones read like DSLD label claims.

\section{Open problems}

\begin{enumerate}
  \item \textbf{Claims belong to products, not ingredients.} A label claim is made about the whole
        product. For single-ingredient products the link is clear. For multi-ingredient products it is
        often unclear which ingredient, if any, a claim is about. The same holds for NHP.
  \item \textbf{Most claims are not about diseases.} iDISK's effect relations point to diseases and
        symptoms, but by a rough keyword estimate about 90\% of DSLD purposes are about normal body
        functions (``supports healthy digestion''), as are 242 of the 269 authorised EU claims. iDISK
        has no concept for these.
  \item \textbf{Positive and negative effects.} Label purposes and EU claims describe benefits, and
        label risks describe harms. Both need to be represented, possibly as one general effect
        relation with a direction, but this is not settled.
  \item \textbf{Claims are not evidence of effect.} A label claim is what a manufacturer states, and a
        non-authorised EU claim is one EFSA rejected. The graph needs to record where a relation comes
        from and how strong its support is, so these are not read as established facts.
  \item \textbf{Matching names across sources.} Label ingredients, NHP ingredients and EU substances
        have to be matched to the same iDISK ingredient. iDISK 2.0 merges entities on ``same CUI or
        same name'', which over-merges, whereas iDISK 1.0 used a stricter process.
  \item \textbf{Data quality.} The EU register has leftover HTML, encoding errors, cells holding
        several claims, and one claim repeated over several rows.
\end{enumerate}

\section{Next steps}

\begin{enumerate}
  \item \textbf{Evaluate the extraction} (see the progress report), so later steps start from
        measured quality.
  \item \textbf{Work out what each claim is about}: first the entities named in the claims (diseases,
        symptoms, body functions), then which ingredient a claim belongs to.
  \item \textbf{Decide how claims enter iDISK}, once the entities are known: which relation types and
        concept types are needed, and how source and verdict are recorded.
  \item \textbf{Bring in the EU register} the same way, as a regulator-checked reference for the
        ingredients it covers.
\end{enumerate}

\end{document}
